\chapter{Introduction}
\label{ch:intro}

Retrieval-augmented generation (RAG) has become the default way to put a large
language model (LLM) in front of a body of knowledge it was not trained on
\cite{lewis2020rag}. A retriever selects a handful of passages relevant to the
user's question from a document collection. The model is then asked to answer
from those passages. The arrangement solves two practical problems at once: the
model can use knowledge newer than its training data, and its answers can be
traced back to sources.

It also creates an attack surface. The document collection is frequently
assembled from sources the operator does not fully control: web pages, shared
drives, ticketing systems, user uploads. Anyone who can place text in it can
place text in front of the model. The literature on this threat has
concentrated on \emph{integrity}: crafting passages that make the model give a
specific wrong answer \cite{zou2025poisonedrag}, or that are retrieved for many
unrelated queries \cite{zhong2023poisoning}. This thesis studies
\emph{availability}. It asks whether an attacker can make a RAG system stop
answering, and whether some models are easier to silence than others.

\section{The question}
\label{sec:question}

The central question is how alignment training affects a RAG system's
susceptibility to denial of service. Here alignment training means the
instruction tuning and preference optimisation that turn a pretrained
model into an assistant \cite{ouyang2022instructgpt,wei2022flan}.

There are good reasons to expect alignment to help. Aligned models follow
instructions, are more selective about what they assert, and are routinely
deployed with an explicit instruction to answer only from the provided context.
There are equally good reasons to expect it to hurt. An instruction to decline
when the evidence is inadequate is also an instruction an attacker can trigger:
make the evidence \emph{look} inadequate, and a faithful model will decline.
A base model, which has never been taught to decline, may simply keep
answering.

The hypothesis tested is the second of these, stated so that it can fail:

\begin{quote}
\textbf{Hypothesis.} Under the same retrieved context and the same attack, an
instruction-tuned model refuses legitimate queries significantly more often than
its base model, and the attack itself, rather than the model's clean refusal
rate, is responsible for the difference.
\end{quote}

The second clause matters. An aligned model that already declines a third of
queries with no attack present will show a higher refusal rate under attack for
that reason alone. A result built on that would be an artefact of the refusal
floor, not a vulnerability. \Cref{sec:stats} describes how the two are
separated.

\section{Why it matters}
\label{sec:significance}

Three groups have a stake in the answer.
\begin{itemize}
  \item \textbf{Deployers} choose which model to put behind a retriever, and
    usually choose the aligned one. If that choice trades misinformation risk
    for availability risk, they should know the exchange rate.
  \item \textbf{Alignment research} treats refusal as the safe failure mode. A
    result showing that refusal can be induced on benign queries by content
    alone, with no unsafe text involved, bears on how over-refusal is measured
    \cite{rottger2024xstest,cui2024orbench}.
  \item \textbf{RAG security} has so far evaluated attacks on one model at a
    time. A matched base-versus-aligned design isolates the training effect
    that single-model evaluations cannot see.
\end{itemize}
\todo{One paragraph of concrete deployment context: enterprise search,
customer support; cost of an unanswered query versus a wrong one.}

\section{Contributions}
\label{sec:contributions}

\begin{enumerate}
  \item \textbf{\bench{}}, a benchmark of 13 denial-of-service attacks in four
    families (semantic jamming, resource exhaustion, retrieval-layer
    manipulation, and reasoning deadlocks). It is run against four matched
    model pairs on two corpora with 1{,}000 queries per run, and each result is
    tested for both a model difference and an attack effect
    (\cref{ch:method}).
  \item \textbf{A validated measurement instrument.} The refusal classifier
    that scores every result is validated against 300 hand-labelled responses
    and an independent annotator. It agrees with that annotator better than
    two humans agree with each other (\cref{sec:classifier}).
  \item \textbf{The alignment paradox, and its mechanism.} Aligned models are
    the more vulnerable side in every true base-versus-aligned run on both
    corpora. The mechanism is context-faithfulness, not safety. A
    prompt-by-training experiment shows that instruction tuning is what makes
    an instruction to decline operative (\cref{ch:results}).
  \item \textbf{A correctness-based reframing.} The same attack causes an
    aligned model to refuse and a base model to answer wrongly. Alignment
    turns silent error into visible refusal (\cref{sec:res-correctness}).
  \item \textbf{A tested defence, and a counter-intuitive negative result.} A
    near-duplicate cluster filter neutralises the templated attacks at no
    measurable clean cost. Cross-encoder reranking, a standard retrieval
    improvement, makes two of the attacks roughly twice as effective
    (\cref{ch:defences}).
\end{enumerate}

Two further results are negative and are reported as such. None of the attacks
is a denial-of-service attack in the compute sense: median latency inflation is
under 10\% (\cref{sec:res-latency}). And one of the two corpora originally used,
HotpotQA, was withdrawn after an audit found that the corpus builder had
discarded half of every query's required evidence (\cref{sec:threats}).

\section{Report outline}

\Cref{ch:background} covers RAG, alignment training and prior attacks on RAG.
\Cref{ch:method} defines the threat model, the benchmark, the metrics and the
statistical protocol. \Cref{ch:results} reports the main results and the
experiments that establish their mechanism. \Cref{ch:defences} tests defences.
\Cref{ch:discussion} assembles the results into one account and states its
limits, and \cref{ch:conclusion} concludes.
